% !PRINTED 202
\weberpage{202}
Diese Bemerkung giebt uns eine wesentliche Begrenzung der Zahlen $\nu$. Es ist nämlich danach:

\weq{13}{2n - 2 = \mu(\nu - 1) + \mu'(\nu' - 1) + \mu''(\nu'' - 1) + \cdots}

oder, wenn wir mit $h$ die Anzahl der Systeme conjugirter Pole bezeichnen, und $n = \mu \nu = \mu' \nu' = \cdots$ setzen:

\weq{14}{2n - 2 = nh - \mu - \mu' - \mu'' - \cdots}

Die Zahlen $\nu, \nu', \nu'' \ldots$ sind mindestens gleich 2, also ist

\[ 1 \leqq \mu \leqq \frac{n}{2}, \quad 1 \leqq \mu' \leqq \frac{n}{2}, \cdots \]

und daher nach (14)

\[ \frac{nh}{2} \leqq 2n - 2 \leqq (n-1)h, \]

oder

\[ 2 \leqq h \leqq 4 - \frac{4}{n}. \]

Es kann also $h$ nur einen der beiden Werthe 2 oder 3 haben, und wir finden fünf Arten, die Gleichung (13) zu befriedigen.

Wenn zunächst $h = 2$ ist, so folgt aus (14)

\[ \mu + \mu' = 2, \quad \mu = \mu' = 1, \quad \nu = \nu' = n. \]

Wir haben also:

I. \gesperrt{Kreistheilungsgruppe} oder \gesperrt{cyklische Gruppe:}

\[ \nu = \nu' = n, \quad n \text{ beliebig.} \]
\[ \mu = \mu' = 1, \]

Ist ferner $h = 3$, so folgt aus (14):

\weq{15}{\mu + \mu' + \mu'' = n + 2.}

Daraus ist zu schliessen, dass mindestens eine der Zahlen $\nu, \nu', \nu''$ gleich 2 sein muss. Denn wären sie alle $\geqq 3$, so wäre

\[ \mu \leqq \frac{n}{3}, \quad \mu' \leqq \frac{n}{3}, \quad \mu'' \leqq \frac{n}{3}, \]

also $\mu + \mu' + \mu'' \leqq n$, was mit (15) im Widerspruch steht.

Ist also $\nu = 2$, $\mu = \dfrac{n}{2}$, so folgt aus (15):

\weq{16}{\mu' + \mu'' = \frac{n}{2} + 2.}

Wir nehmen zunächst an, dass auch noch $\nu' = 2$ sei. Setzen wir dann $\nu'' = m$, so folgt aus (16):

\[ \mu'' = 2, \quad n = 2m, \]
%CONT>>
